At a serial CPU proposal budget of 5\,ms, multisite exact-composition drifting recovered 43.75\% of targets, the one-step flow 47.66\%, the 16-step flow 11.72\%, and retrieval 57.55\% (Fig.~\ref{fig:time}). Given five proposals each, the 16-step flow overtook direct drifting (Table~\ref{tab:recovery}). The four-step flow on two-site cells illustrates how a smaller per-step network changes the comparison. It recovered 82.72\% in five and cost 0.714\,ms per proposal on CPU. The budget estimates use each seed's own measured median and stop at the observed five draws.

We also replayed the released 131-million-parameter Cartesian uFlowCSP exponential-moving-average checkpoint with a single Euler update and no subsequent relaxation or ranking \cite{uflow}. All 189 two-site and 128 larger-cohort IDs were in its published test split. It recovered 174/189 in five draws versus 82.01\% for our two-site model, a paired gap of 10.05 percentage points (95\% case-bootstrap interval $[5.64,14.64]$). At five draws per target, its two-site geometric energy score was 0.192, versus 1.219 for exact drifting and 0.510 for retrieval. On the larger cohort it recovered 103/128 in five draws, including 87/112 with three or more sites. Its advantage over exact drifting was 36.72 percentage points ($[26.82,45.83]$). The authors used the substantially larger MP-20 training corpus. We timed that checkpoint separately with batches of 12 on three CPU threads.
